\chapter{Data Strategy and Costs}\label{ch:fm:data-strategy-and-costs}

In February 2024 the UK's Financial Conduct Authority published the findings of its study of the wholesale data market. It found no evidence that
firms could not get the data they needed, and evidence of market power in all three markets it studied: benchmarks, credit ratings data and the
market data vendors that redistribute trading data. Users, it concluded, may be paying higher prices than competition would allow; licensing was
complex, many users had to hold licences both from a trading venue and from the vendor that delivered its data, and the complexity itself cost money,
down to the compliance teams firms employ to follow their licences. A trading firm's data bill is set less by how much data it uses than by how it
is counted.

\section{What a firm pays for data}

\begin{definition}[Data budget]\label{def:fm:data-strategy-and-costs:budget}
A \emph{data budget}\index{data budget} is a firm's planned annual spending on data, by product and use: exchange market data (display,
non-display and feeds), vendor terminals and platforms, reference and corporate-action data, news, and alternative datasets, with the licences,
counts and allocations behind each line.
\end{definition}

The chapter's firm (\cref{tab:fm:data-strategy-and-costs:inventory}) spends \$2.07 million a year on data: \$0.56 million on exchange data, \$0.75
million on thirty vendor terminals, \$0.46 million on reference data and news, and \$0.30 million on an alternative dataset. Terminals are the largest
line, 36\% of the total, and the one most directly driven by headcount. The fee levels are illustrative inputs; their structure is what every data
contract shares: something charged per user, per device, per use, or as a flat fee, with a cap or an enterprise alternative.

\begin{table}[ht]
\centering\small\setlength{\tabcolsep}{4pt}
\begin{tabular}{llrr}
\toprule
product & charged & usage & \$ thousand a year \\
\midrule
exchange real-time display & 0.6 per user & 40 users & 24 \\
exchange non-display (automated) & flat & -- & 180 \\
direct feeds & flat, six feeds & -- & 360 \\
vendor terminals & 25 per user & 30 users & 750 \\
reference data & 8 per user or 400 & 60 users & 400 \\
news analytics & 15 per device & 4 devices & 60 \\
alternative dataset & flat & -- & 300 \\
\midrule
total & & & 2\,074 \\
\bottomrule
\end{tabular}
\caption{The chapter's data inventory, each product at its cheaper licence (illustrative fee levels). Data: \texttt{fm\_data}.}\label{tab:fm:data-strategy-and-costs:inventory}
\end{table}

\begin{omfigure}
\begin{tikzpicture}
\begin{axis}[width=10.5cm,height=4.8cm,xbar,bar width=10pt,xmin=0,xmax=900,xlabel={\small \$ thousand a year},ytick={0,1,2,3},
  yticklabels={market data,terminals,reference data,alternative data},y dir=reverse,ymin=-0.6,ymax=3.6,
  tick label style={font=\footnotesize},grid=major,grid style={gray!20},table/col sep=comma,nodes near coords,
  nodes near coords style={font=\scriptsize}]
\addplot[fill=omDef!55,draw=omDef] table[y=pos,x=cost]{figdata/desk/22-data-strategy-and-costs/budget.csv};
\end{axis}
\end{tikzpicture}

\omcaption{The \omterm{def:fm:data-strategy-and-costs:budget}{data budget} by category, \$2.07 million a year: terminals, the headcount-driven line, are the largest. Data:
\texttt{fm\_data.the\_budget}.}\label{fig:fm:data-strategy-and-costs:budget}
\end{omfigure}

\section{Licences and their units of count}

Book 15 defined the unit of count, the entitlement system and the usage report: the machinery by which a firm knows who and what consumes each
dataset. The licences they serve come in a few shapes.

\begin{definition}[Enterprise, derived-data and redistribution licences]\label{def:fm:data-strategy-and-costs:licences}
An \emph{enterprise licence}\index{enterprise licence} permits unlimited use of a product within the firm (or a defined part of it) for a fixed fee,
replacing per-user or per-device counts. A \emph{derived-data licence}\index{derived-data licence} permits the firm to create and use, and
sometimes distribute, data calculated from the licensed data, such as prices, indices or signals. A \emph{redistribution
licence}\index{redistribution licence} permits the firm to pass the data, or data derived from it, to third parties such as clients, under
conditions and for fees.
\end{definition}

The enterprise decision is a break-even. The reference data costs \$8\,000 per user a year or \$400\,000 for the enterprise: the \omterm{def:fm:data-strategy-and-costs:licences}{enterprise licence}
wins from 50 users (\cref{fig:fm:data-strategy-and-costs:enterprise}). At the firm's sixty users it saves \$80\,000 a year, and it removes the
counting, and with it the audit exposure, for that product. Below fifty users per-user fees are cheaper, unless headcount is expected to grow past
the break-even within the contract.

\begin{omfigure}
\begin{tikzpicture}
\begin{axis}[width=10.5cm,height=5cm,xlabel={\small users},ylabel={\small \$ thousand a year},xmin=0,xmax=100,ymin=0,ymax=820,
  tick label style={font=\footnotesize},grid=major,grid style={gray!20},table/col sep=comma,
  legend cell align=left,legend style={font=\scriptsize,at={(0.5,-0.3)},anchor=north,/tikz/every even column/.append style={column sep=8pt}},legend columns=2]
\addplot[omDef,thick] table[x=users,y=per_user]{figdata/desk/22-data-strategy-and-costs/enterprise.csv};
\addplot[omProp,thick,dashed] table[x=users,y=enterprise]{figdata/desk/22-data-strategy-and-costs/enterprise.csv};
\draw[gray,dotted] (axis cs:50,0) -- (axis cs:50,820);
\draw[gray,dotted] (axis cs:60,0) -- (axis cs:60,820);
\legend{per-user fees,enterprise licence}
\end{axis}
\end{tikzpicture}

\omcaption{Per-user fees against an \omterm{def:fm:data-strategy-and-costs:licences}{enterprise licence} for the reference data: the lines cross at 50 users; the firm has 60. Data:
\texttt{fm\_data.reference\_curve}.}\label{fig:fm:data-strategy-and-costs:enterprise}
\end{omfigure}

\omcode{code/firm/databudget/firm_databudget.py}{36}{67}{A product's annual cost under its fee rules, the budget at the cheaper licence, the
enterprise break-even and the exposure of an audit.}\label{lst:fm:data-strategy-and-costs:budget}

\section{Audits and back-billing}

\begin{definition}[Back-billing]\label{def:fm:data-strategy-and-costs:backbill}
\emph{Back-billing}\index{back-billing} is the charge a data provider raises after an audit for usage that was not reported or paid for in past
periods, usually with interest and sometimes with a penalty, under the audit clause of its licence.
\end{definition}

\begin{proposition}[What an audit finds]\label{prop:fm:data-strategy-and-costs:audit}
If a firm reported fees $F$ a year while a share $s$ of its true usage went unreported for $M$ months, the unbilled fees are $u=Fs/(12(1-s))$ a
month and the \omterm{def:fm:data-strategy-and-costs:backbill}{back-billing} with monthly interest $i$ from each month to the audit is
\[
u\sum_{k=0}^{M-1}(1+i)^k=u\,\frac{(1+i)^M-1}{i}.
\]
\end{proposition}

\begin{proof}
True usage is reported usage divided by $1-s$, so the unreported fees are $F\,s/(1-s)$ a year. The month $m$'s unbilled amount accrues interest for
$M-m$ months; summing the geometric series gives the formula.
\end{proof}

The firm's per-user fees, display and terminals, are \$774\,000 a year. An audit that finds 15\% of true users unreported over three years bills
\$410\,000 of fees and, at 1\% a month, \$81\,000 of interest: \$490\,000, 63\% of a year's per-user fees, due at once
(\cref{fig:fm:data-strategy-and-costs:audit}). The exposure grows faster than the share: 5\% under-reporting costs \$146\,000, 30\% costs
\$1.19 million. The defence is the entitlement system of Book 15: counts taken from the systems that grant access, not from what users declare.

\begin{omfigure}
\begin{tikzpicture}
\begin{axis}[width=10.5cm,height=5cm,ybar stacked,bar width=16pt,xlabel={\small share of true users unreported (\%)},ylabel={\small \$ thousand},
  ymin=0,ymax=1300,xtick={5,10,15,20,25,30},tick label style={font=\footnotesize},grid=major,grid style={gray!20},table/col sep=comma,
  legend cell align=left,legend style={font=\scriptsize,at={(0.5,-0.3)},anchor=north,/tikz/every even column/.append style={column sep=8pt}},legend columns=2]
\addplot[fill=omDef!55,draw=omDef] table[x=under,y=principal]{figdata/desk/22-data-strategy-and-costs/audit.csv};
\addplot[fill=omProp!55,draw=omProp] table[x=under,y=interest]{figdata/desk/22-data-strategy-and-costs/audit.csv};
\legend{unbilled fees,interest at 1\% a month}
\end{axis}
\end{tikzpicture}

\omcaption{\omterm{def:fm:data-strategy-and-costs:backbill}{Back-billing} after an audit of three years of per-user fees (\$774\,000 a year reported), by the share of true users that went
unreported. Data: \texttt{fm\_data.audit}.}\label{fig:fm:data-strategy-and-costs:audit}
\end{omfigure}

\section{Tutorial: the audit letter}\label{tut:fm:data-strategy-and-costs:tut}

\begin{tutorial}
\textbf{Goal.} Price a firm's data inventory, find where an \omterm{def:fm:data-strategy-and-costs:licences}{enterprise licence} wins, and estimate the exposure of an audit that finds users
unreported. \textbf{End state:} the budget by category (\cref{fig:fm:data-strategy-and-costs:budget}), the enterprise break-even
(\cref{fig:fm:data-strategy-and-costs:enterprise}) and the audit exposure (\cref{fig:fm:data-strategy-and-costs:audit}).
\begin{tutsteps}
\item \textbf{The inventory.} \texttt{fm\_data.PRODUCTS} and \texttt{USAGE} hold seven products with their fee rules and counts;
\texttt{firm.databudget.budget} prices each at its cheaper licence.
\item \textbf{The break-even.} \texttt{enterprise\_breakeven} gives the headcount from which the \omterm{def:fm:data-strategy-and-costs:licences}{enterprise licence} is cheaper;
\texttt{fm\_data.reference\_curve} draws both costs against users.
\item \textbf{The audit.} \texttt{fm\_data.audit(under, years, monthly\_rate)} applies \cref{prop:fm:data-strategy-and-costs:audit} to the
per-user fees.
\item \textbf{The allocation.} \texttt{fm\_data.allocation()} splits the budget over four strategies by their measured use;
\texttt{fm\_data.alt\_value()} prices the alternative dataset with Book 7's vendor-trial arithmetic.
\end{tutsteps}
\end{tutorial}

\textbf{What to change next.} Let headcount grow 20\% a year and find when the terminals would justify an enterprise negotiation; add a
redistribution fee for a client-facing quote page.

\section{Negotiating with exchanges and vendors}

\begin{dated}{2026-09}{Rules on market-data pricing}\label{dat:fm:data-strategy-and-costs:rules}
\textbf{European Union}: MiFIR (Regulation (EU) No 600/2014) article 13, as adopted, requires trading venues to make pre- and post-trade data
available on a reasonable commercial basis, with non-discriminatory access, and free of charge fifteen minutes after publication. \textbf{United
States}: the SEC's market data infrastructure rules provide for competing consolidators of NMS market data, registered under Rule 614 (17 CFR
242.614, adopted in 2021). \textbf{United Kingdom}: the FCA's wholesale data market study (February 2024) found evidence of market power among
benchmark providers, credit rating agencies and market data vendors, and said it would pursue its findings through regulatory reform and its
competition powers.
\end{dated}

Negotiation starts from the inventory and the counts. A firm that knows exactly who uses what can move whole categories to enterprise terms, drop
unused entitlements before renewal, and resist price rises with a credible plan to switch where a substitute exists. It has less leverage where the
data comes from one venue or one administrator, which is the market power the FCA described; there the levers are the licence terms rather than
the price: definitions of non-display use, derived-data rights, audit periods and interest, and caps on renewal increases.

\section{Data as a strategic asset}

An alternative dataset is bought for the return it adds. Book 7's vendor-trial harness prices it: the annual fee at which the dataset's expected
contribution, decaying as others buy it, just pays for itself. For the chapter's dataset, feeding a strategy on \$50 million that earns 2\% a year
from it before trading costs of 0.5\%, with an edge that halves every two years, the break-even fee over a three-year contract is \$369\,000 a year
(\texttt{fm\_data.alt\_value}): at \$300\,000 the dataset is worth buying, with little margin for error in the decay.

Data costs are allocated to strategies by use: on the chapter's weights, market making carries \$933\,000 of the budget, statistical arbitrage
\$622\,000, event-driven trading \$311\,000 and macro \$207\,000. A strategy whose data costs exceed what it earns from them should not be
carrying them, which is only visible once they are allocated.

\begin{method}[Running a data budget]\label{met:fm:data-strategy-and-costs:run}
\begin{enumerate}
\item Keep an inventory of every data product: provider, licence, unit of count, current counts, renewal date, owner.
\item Take counts from the entitlement system, reconcile them monthly with providers' invoices, and remove unused access before each renewal.
\item Compare per-unit and enterprise terms at forecast headcount; price the audit exposure of every counted licence.
\item Price alternative data by its expected contribution before and after a trial, and stop what does not pay.
\item Allocate the budget to strategies by use and review it with their results.
\end{enumerate}
\end{method}

\section{Build: the data budget}\label{bld:fm:data-strategy-and-costs:databudget}

\begin{build}
\bfield{Purpose} Data products and fee rules as data, the budget, the enterprise break-even, audit exposure, allocation to strategies and the value
of an alternative dataset.
\bfield{Interface} \texttt{firm.databudget}: \texttt{Product}, \texttt{Usage}, \texttt{annual\_cost}, \texttt{budget}, \texttt{enterprise\_breakeven},
\texttt{audit\_exposure}, \texttt{allocate}, \texttt{alt\_breakeven} (wrapping \texttt{firm.vendoreval.breakeven}).
\bfield{Rules} Fee levels are inputs; each product is budgeted at its cheaper licence; \omterm{def:fm:data-strategy-and-costs:backbill}{back-billing} accrues interest from each month to the audit.
\bfield{Acceptance tests} \texttt{code/firm/databudget/tests/}: costs and the budget on hand numbers; the break-even; the audit's geometric sum; the
wrapper.
\bfield{Stretch} Usage records from Book 15's entitlement system; renewal dates and notice periods; derived-data and redistribution fees.
\end{build}

\begin{omsources}
\item Financial Conduct Authority, \emph{Wholesale Data Market Study}, Market Study MS23/1.5, February 2024.
\item Regulation (EU) No 600/2014, article 13; 17 CFR 242.614.
\end{omsources}

\section{Exercises}

\begin{exercise}[$\star$]\label{exo:fm:data-strategy-and-costs:1}
At what headcount does the reference data's \omterm{def:fm:data-strategy-and-costs:licences}{enterprise licence} win, and what does it save at sixty users?
\end{exercise}

\begin{exercise}[$\star$]\label{exo:fm:data-strategy-and-costs:2}
What did the FCA's study find about access to data and about prices?
\end{exercise}

\begin{exercise}[$\star$]\label{exo:fm:data-strategy-and-costs:3}
Without interest, what does an audit bill for 15\% of true users unreported over three years on \$774\,000 a year of per-user fees?
\end{exercise}

\begin{exercise}[$\star\star$]\label{exo:fm:data-strategy-and-costs:4}
Prove the closed form of \cref{prop:fm:data-strategy-and-costs:audit} and compute the interest at 1\% a month.
\end{exercise}

\begin{exercise}[$\star\star$]\label{exo:fm:data-strategy-and-costs:5}
Explain the difference between a \omterm{def:fm:data-strategy-and-costs:licences}{derived-data licence} and a \omterm{def:fm:data-strategy-and-costs:licences}{redistribution licence} with an example from a trading firm.
\end{exercise}

\begin{exercise}[$\star\star$]\label{exo:fm:data-strategy-and-costs:6}
Why does the audit exposure grow faster than the unreported share?
\end{exercise}

\begin{exercise}[$\star\star\star$]\label{exo:fm:data-strategy-and-costs:7}
\emph{Coding.} With a half-life of one year instead of two, what is the alternative dataset's break-even fee? Should the firm buy at \$300\,000?
\end{exercise}

\begin{exercise}[$\star\star\star$]\label{exo:fm:data-strategy-and-costs:8}
\emph{Find the flaw.} ``Our users tell us which terminals they use each quarter, so our counts are right.''
\end{exercise}

\section{Problem: The Audit Letter}

\begin{problem}[{Weekend problem --- the audit letter}]\label{pb:fm:data-strategy-and-costs:1}
A chief operating officer receives an exchange's notice of audit, and must estimate the exposure and redesign how the firm buys data.

\medskip\noindent\textbf{Part I --- The market.}
\begin{enumerate}
\item What did the FCA find in 2024, and why does it matter for a buyer?
\item What does MiFIR article 13 require of trading venues?
\item Define a \omterm{def:fm:data-strategy-and-costs:budget}{data budget} and list its categories.
\item Give the firm's budget by category.
\end{enumerate}

\medskip\noindent\textbf{Part II --- Licences.}
\begin{enumerate}[resume]
\item Define enterprise, derived-data and \omterm{def:fm:data-strategy-and-costs:licences}{redistribution licences}.
\item Compute the reference data's enterprise break-even.
\item What else, besides price, does an \omterm{def:fm:data-strategy-and-costs:licences}{enterprise licence} change?
\item Which terms would you negotiate where the provider has market power?
\end{enumerate}

\medskip\noindent\textbf{Part III --- The audit.}
\begin{enumerate}[resume]
\item Define \omterm{def:fm:data-strategy-and-costs:backbill}{back-billing}.
\item State and prove \cref{prop:fm:data-strategy-and-costs:audit}.
\item Give the exposure for 15\% under-reporting over three years, and for 5\% and 30\%.
\item How would you respond to the notice?
\item How does an entitlement system reduce the exposure?
\end{enumerate}

\medskip\noindent\textbf{Part IV --- Strategy.}
\begin{enumerate}[resume]
\item How would you price the alternative dataset, and what is its break-even fee?
\item Allocate the budget to strategies.
\item What does the allocation tell the firm?
\item What would you stop buying, and how would you decide?
\item How often should the inventory be reconciled?
\item State the \emph{named result}: the headcount at which an \omterm{def:fm:data-strategy-and-costs:licences}{enterprise licence} beats per-user fees, and the \omterm{def:fm:data-strategy-and-costs:backbill}{back-billing} exposure of an audit that
finds 15 per cent under-reporting over three years.
\item In two sentences, write the data policy.
\end{enumerate}
\end{problem}

\section{Interview questions}

\begin{interviewq}[$\star$ \iqroles{developer}]\label{iq:fm:data-strategy-and-costs:1}
What is a non-display use of market data, and why does it matter to a trading firm?
\end{interviewq}

\begin{interviewq}[$\star$ \iqroles{trader}]\label{iq:fm:data-strategy-and-costs:2}
Why might a firm pay for an \omterm{def:fm:data-strategy-and-costs:licences}{enterprise licence} it does not fully use?
\end{interviewq}

\begin{interviewq}[$\star\star$ \iqroles{researcher}]\label{iq:fm:data-strategy-and-costs:3}
How would you decide whether an alternative dataset is worth its price?
\end{interviewq}

\begin{interviewq}[$\star\star$ \iqroles{developer}]\label{iq:fm:data-strategy-and-costs:4}
Design a system that keeps market-data counts audit-ready.
\end{interviewq}

\begin{interviewq}[$\star\star$ \iqroles{risk}]\label{iq:fm:data-strategy-and-costs:5}
An exchange audit finds under-reporting. What is the firm's exposure, and what do you do next?
\end{interviewq}

\begin{interviewq}[$\star\star\star$ \iqroles{researcher, trader}]\label{iq:fm:data-strategy-and-costs:6}
How should data costs be allocated across strategies, and what behaviour does the allocation create?
\end{interviewq}
